% Options for packages loaded elsewhere
\PassOptionsToPackage{unicode}{hyperref}
\PassOptionsToPackage{hyphens}{url}
%
\documentclass[
]{article}
\usepackage{amsmath,amssymb}
\usepackage{iftex}
\ifPDFTeX
  \usepackage[T1]{fontenc}
  \usepackage[utf8]{inputenc}
  \usepackage{textcomp} % provide euro and other symbols
\else % if luatex or xetex
  \usepackage{unicode-math} % this also loads fontspec
  \defaultfontfeatures{Scale=MatchLowercase}
  \defaultfontfeatures[\rmfamily]{Ligatures=TeX,Scale=1}
\fi
\usepackage{lmodern}
\ifPDFTeX\else
  % xetex/luatex font selection
\fi
% Use upquote if available, for straight quotes in verbatim environments
\IfFileExists{upquote.sty}{\usepackage{upquote}}{}
\IfFileExists{microtype.sty}{% use microtype if available
  \usepackage[]{microtype}
  \UseMicrotypeSet[protrusion]{basicmath} % disable protrusion for tt fonts
}{}
\makeatletter
\@ifundefined{KOMAClassName}{% if non-KOMA class
  \IfFileExists{parskip.sty}{%
    \usepackage{parskip}
  }{% else
    \setlength{\parindent}{0pt}
    \setlength{\parskip}{6pt plus 2pt minus 1pt}}
}{% if KOMA class
  \KOMAoptions{parskip=half}}
\makeatother
\usepackage{xcolor}
\usepackage[margin=1in]{geometry}
\usepackage{color}
\usepackage{fancyvrb}
\newcommand{\VerbBar}{|}
\newcommand{\VERB}{\Verb[commandchars=\\\{\}]}
\DefineVerbatimEnvironment{Highlighting}{Verbatim}{commandchars=\\\{\}}
% Add ',fontsize=\small' for more characters per line
\usepackage{framed}
\definecolor{shadecolor}{RGB}{248,248,248}
\newenvironment{Shaded}{\begin{snugshade}}{\end{snugshade}}
\newcommand{\AlertTok}[1]{\textcolor[rgb]{0.94,0.16,0.16}{#1}}
\newcommand{\AnnotationTok}[1]{\textcolor[rgb]{0.56,0.35,0.01}{\textbf{\textit{#1}}}}
\newcommand{\AttributeTok}[1]{\textcolor[rgb]{0.13,0.29,0.53}{#1}}
\newcommand{\BaseNTok}[1]{\textcolor[rgb]{0.00,0.00,0.81}{#1}}
\newcommand{\BuiltInTok}[1]{#1}
\newcommand{\CharTok}[1]{\textcolor[rgb]{0.31,0.60,0.02}{#1}}
\newcommand{\CommentTok}[1]{\textcolor[rgb]{0.56,0.35,0.01}{\textit{#1}}}
\newcommand{\CommentVarTok}[1]{\textcolor[rgb]{0.56,0.35,0.01}{\textbf{\textit{#1}}}}
\newcommand{\ConstantTok}[1]{\textcolor[rgb]{0.56,0.35,0.01}{#1}}
\newcommand{\ControlFlowTok}[1]{\textcolor[rgb]{0.13,0.29,0.53}{\textbf{#1}}}
\newcommand{\DataTypeTok}[1]{\textcolor[rgb]{0.13,0.29,0.53}{#1}}
\newcommand{\DecValTok}[1]{\textcolor[rgb]{0.00,0.00,0.81}{#1}}
\newcommand{\DocumentationTok}[1]{\textcolor[rgb]{0.56,0.35,0.01}{\textbf{\textit{#1}}}}
\newcommand{\ErrorTok}[1]{\textcolor[rgb]{0.64,0.00,0.00}{\textbf{#1}}}
\newcommand{\ExtensionTok}[1]{#1}
\newcommand{\FloatTok}[1]{\textcolor[rgb]{0.00,0.00,0.81}{#1}}
\newcommand{\FunctionTok}[1]{\textcolor[rgb]{0.13,0.29,0.53}{\textbf{#1}}}
\newcommand{\ImportTok}[1]{#1}
\newcommand{\InformationTok}[1]{\textcolor[rgb]{0.56,0.35,0.01}{\textbf{\textit{#1}}}}
\newcommand{\KeywordTok}[1]{\textcolor[rgb]{0.13,0.29,0.53}{\textbf{#1}}}
\newcommand{\NormalTok}[1]{#1}
\newcommand{\OperatorTok}[1]{\textcolor[rgb]{0.81,0.36,0.00}{\textbf{#1}}}
\newcommand{\OtherTok}[1]{\textcolor[rgb]{0.56,0.35,0.01}{#1}}
\newcommand{\PreprocessorTok}[1]{\textcolor[rgb]{0.56,0.35,0.01}{\textit{#1}}}
\newcommand{\RegionMarkerTok}[1]{#1}
\newcommand{\SpecialCharTok}[1]{\textcolor[rgb]{0.81,0.36,0.00}{\textbf{#1}}}
\newcommand{\SpecialStringTok}[1]{\textcolor[rgb]{0.31,0.60,0.02}{#1}}
\newcommand{\StringTok}[1]{\textcolor[rgb]{0.31,0.60,0.02}{#1}}
\newcommand{\VariableTok}[1]{\textcolor[rgb]{0.00,0.00,0.00}{#1}}
\newcommand{\VerbatimStringTok}[1]{\textcolor[rgb]{0.31,0.60,0.02}{#1}}
\newcommand{\WarningTok}[1]{\textcolor[rgb]{0.56,0.35,0.01}{\textbf{\textit{#1}}}}
\usepackage{graphicx}
\makeatletter
\newsavebox\pandoc@box
\newcommand*\pandocbounded[1]{% scales image to fit in text height/width
  \sbox\pandoc@box{#1}%
  \Gscale@div\@tempa{\textheight}{\dimexpr\ht\pandoc@box+\dp\pandoc@box\relax}%
  \Gscale@div\@tempb{\linewidth}{\wd\pandoc@box}%
  \ifdim\@tempb\p@<\@tempa\p@\let\@tempa\@tempb\fi% select the smaller of both
  \ifdim\@tempa\p@<\p@\scalebox{\@tempa}{\usebox\pandoc@box}%
  \else\usebox{\pandoc@box}%
  \fi%
}
% Set default figure placement to htbp
\def\fps@figure{htbp}
\makeatother
\setlength{\emergencystretch}{3em} % prevent overfull lines
\providecommand{\tightlist}{%
  \setlength{\itemsep}{0pt}\setlength{\parskip}{0pt}}
\setcounter{secnumdepth}{-\maxdimen} % remove section numbering
\usepackage{bookmark}
\IfFileExists{xurl.sty}{\usepackage{xurl}}{} % add URL line breaks if available
\urlstyle{same}
\hypersetup{
  pdftitle={ANA515 Project Part 1 - R Script},
  pdfauthor={Anushree Das},
  hidelinks,
  pdfcreator={LaTeX via pandoc}}

\title{ANA515 Project Part 1 - R Script}
\author{Anushree Das}
\date{}

\begin{document}
\maketitle

\begin{Shaded}
\begin{Highlighting}[]
\CommentTok{\#Load necessary libraries}
\FunctionTok{library}\NormalTok{(dplyr)}
\FunctionTok{library}\NormalTok{(lubridate)}
\FunctionTok{library}\NormalTok{(tidyverse)}
\end{Highlighting}
\end{Shaded}

\section{1. EXTRACT}\label{extract}

\begin{Shaded}
\begin{Highlighting}[]
\CommentTok{\# Load dataset from csv}
\NormalTok{users\_csv  }\OtherTok{\textless{}{-}} \FunctionTok{read.csv}\NormalTok{(}\StringTok{"credit\_card\_transactions/sd254\_users.csv"}\NormalTok{)}
\NormalTok{cards\_csv }\OtherTok{\textless{}{-}} \FunctionTok{read.csv}\NormalTok{(}\StringTok{"credit\_card\_transactions/sd254\_cards.csv"}\NormalTok{)}
\NormalTok{transactions\_csv }\OtherTok{\textless{}{-}} \FunctionTok{read.csv}\NormalTok{(}\StringTok{"credit\_card\_transactions/credit\_card\_transactions{-}ibm\_v2.csv"}\NormalTok{)}
\end{Highlighting}
\end{Shaded}

\section{2. TRANFORM}\label{tranform}

\subsection{2.1 Modify users table}\label{modify-users-table}

\subsubsection{2.1.a Generate UserId column for users
table}\label{a-generate-userid-column-for-users-table}

\subsubsection{2.1.b Rename column names}\label{b-rename-column-names}

\begin{Shaded}
\begin{Highlighting}[]
\NormalTok{users }\OtherTok{\textless{}{-}} \FunctionTok{data.frame}\NormalTok{(}\AttributeTok{UserId =} \DecValTok{1}\SpecialCharTok{:}\FunctionTok{nrow}\NormalTok{(users\_csv), }
                        \AttributeTok{FullName =}\NormalTok{ users\_csv}\SpecialCharTok{$}\NormalTok{Person, }
                        \AttributeTok{CurrentAge =}\NormalTok{ users\_csv}\SpecialCharTok{$}\NormalTok{Current.Age,}
                        \AttributeTok{RetirementAge =}\NormalTok{ users\_csv}\SpecialCharTok{$}\NormalTok{Retirement.Age,}
                        \AttributeTok{BirthYear =}\NormalTok{ users\_csv}\SpecialCharTok{$}\NormalTok{Birth.Year, }
                        \AttributeTok{BirthMonth =}\NormalTok{ users\_csv}\SpecialCharTok{$}\NormalTok{Birth.Month, }
                        \AttributeTok{Gender=}\NormalTok{ users\_csv}\SpecialCharTok{$}\NormalTok{Gender, }
                        \AttributeTok{Address1 =}\NormalTok{ users\_csv}\SpecialCharTok{$}\NormalTok{Address, }
                        \AttributeTok{Address2 =}\NormalTok{ users\_csv}\SpecialCharTok{$}\NormalTok{Apartment, }
                        \AttributeTok{City =}\NormalTok{ users\_csv}\SpecialCharTok{$}\NormalTok{City, }
                        \AttributeTok{State =}\NormalTok{ users\_csv}\SpecialCharTok{$}\NormalTok{State, }
                        \AttributeTok{Zipcode =}\NormalTok{ users\_csv}\SpecialCharTok{$}\NormalTok{Zipcode, }
                        \AttributeTok{Latitude =}\NormalTok{ users\_csv}\SpecialCharTok{$}\NormalTok{Latitude, }
                        \AttributeTok{Longitude =}\NormalTok{ users\_csv}\SpecialCharTok{$}\NormalTok{Longitude, }
                        \AttributeTok{PerCapitaIncomeZipcode =}\NormalTok{ users\_csv}\SpecialCharTok{$}\NormalTok{Per.Capita.Income...Zipcode, }
                        \AttributeTok{YearlyIncome =}\NormalTok{ users\_csv}\SpecialCharTok{$}\NormalTok{Yearly.Income...Person, }
                        \AttributeTok{TotalDebt =}\NormalTok{ users\_csv}\SpecialCharTok{$}\NormalTok{Total.Debt, }
                        \AttributeTok{FICOScore =}\NormalTok{ users\_csv}\SpecialCharTok{$}\NormalTok{FICO.Score, }
                        \AttributeTok{NumCreditCards =}\NormalTok{ users\_csv}\SpecialCharTok{$}\NormalTok{Num.Credit.Cards}
\NormalTok{)}

\FunctionTok{head}\NormalTok{(users, }\DecValTok{3}\NormalTok{)}
\end{Highlighting}
\end{Shaded}

\begin{verbatim}
##   UserId       FullName CurrentAge RetirementAge BirthYear BirthMonth Gender
## 1      1 Hazel Robinson         53            66      1966         11 Female
## 2      2     Sasha Sadr         53            68      1966         12 Female
## 3      3     Saanvi Lee         81            67      1938         11 Female
##                 Address1 Address2        City State Zipcode Latitude Longitude
## 1          462 Rose Lane       NA    La Verne    CA   91750    34.15   -117.76
## 2 3606 Federal Boulevard       NA Little Neck    NY   11363    40.76    -73.74
## 3        766 Third Drive       NA West Covina    CA   91792    34.02   -117.89
##   PerCapitaIncomeZipcode YearlyIncome TotalDebt FICOScore NumCreditCards
## 1                 $29278       $59696   $127613       787              5
## 2                 $37891       $77254   $191349       701              5
## 3                 $22681       $33483      $196       698              5
\end{verbatim}

\subsubsection{2.1.c Convert Gender to categorical data
type}\label{c-convert-gender-to-categorical-data-type}

\begin{Shaded}
\begin{Highlighting}[]
\NormalTok{users}\SpecialCharTok{$}\NormalTok{Gender }\OtherTok{\textless{}{-}} \FunctionTok{factor}\NormalTok{(users}\SpecialCharTok{$}\NormalTok{Gender)}
\end{Highlighting}
\end{Shaded}

\subsubsection{2.1.d Convert columns - Address2, Zipcode to Text data
type}\label{d-convert-columns---address2-zipcode-to-text-data-type}

Zipcode and Address2 should be converted to Text data type because they
are not used for quantitative statistics

\begin{Shaded}
\begin{Highlighting}[]
\NormalTok{users}\SpecialCharTok{$}\NormalTok{Address2 }\OtherTok{\textless{}{-}} \FunctionTok{as.character}\NormalTok{(users}\SpecialCharTok{$}\NormalTok{Address2)}
\NormalTok{users}\SpecialCharTok{$}\NormalTok{Zipcode }\OtherTok{\textless{}{-}} \FunctionTok{as.character}\NormalTok{(users}\SpecialCharTok{$}\NormalTok{Zipcode)}
\end{Highlighting}
\end{Shaded}

\subsubsection{2.1.e Convert columns - PerCapitaIncomeZipcode,
YearlyIncome, TotalDebt to
numeric}\label{e-convert-columns---percapitaincomezipcode-yearlyincome-totaldebt-to-numeric}

\begin{Shaded}
\begin{Highlighting}[]
\CommentTok{\# Remove $ character from the column values and convert to numeric data type}
\NormalTok{users}\SpecialCharTok{$}\NormalTok{PerCapitaIncomeZipcode }\OtherTok{\textless{}{-}} \FunctionTok{as.numeric}\NormalTok{(}\FunctionTok{sub}\NormalTok{(}\StringTok{\textquotesingle{}.\textquotesingle{}}\NormalTok{, }\StringTok{\textquotesingle{}\textquotesingle{}}\NormalTok{, users}\SpecialCharTok{$}\NormalTok{PerCapitaIncomeZipcode))}
\NormalTok{users}\SpecialCharTok{$}\NormalTok{YearlyIncome }\OtherTok{\textless{}{-}} \FunctionTok{as.numeric}\NormalTok{(}\FunctionTok{sub}\NormalTok{(}\StringTok{\textquotesingle{}.\textquotesingle{}}\NormalTok{, }\StringTok{\textquotesingle{}\textquotesingle{}}\NormalTok{, users}\SpecialCharTok{$}\NormalTok{YearlyIncome))}
\NormalTok{users}\SpecialCharTok{$}\NormalTok{TotalDebt }\OtherTok{\textless{}{-}} \FunctionTok{as.numeric}\NormalTok{(}\FunctionTok{sub}\NormalTok{(}\StringTok{\textquotesingle{}.\textquotesingle{}}\NormalTok{, }\StringTok{\textquotesingle{}\textquotesingle{}}\NormalTok{, users}\SpecialCharTok{$}\NormalTok{TotalDebt))}
\end{Highlighting}
\end{Shaded}

\subsection{2.2 Modify cards tables}\label{modify-cards-tables}

\subsubsection{2.2.a Remove duplicates}\label{a-remove-duplicates}

\begin{Shaded}
\begin{Highlighting}[]
\NormalTok{cards\_csv }\OtherTok{\textless{}{-}}\NormalTok{ cards\_csv }\SpecialCharTok{\%\textgreater{}\%} \FunctionTok{distinct}\NormalTok{() }
\end{Highlighting}
\end{Shaded}

\subsubsection{2.2.b Combine user id and card index to generate
CardId}\label{b-combine-user-id-and-card-index-to-generate-cardid}

\subsubsection{2.2.c Rename column names}\label{c-rename-column-names}

\begin{Shaded}
\begin{Highlighting}[]
\NormalTok{cards }\OtherTok{\textless{}{-}} \FunctionTok{data.frame}\NormalTok{(}\AttributeTok{CardId =} \FunctionTok{paste}\NormalTok{(cards\_csv}\SpecialCharTok{$}\NormalTok{User, cards\_csv}\SpecialCharTok{$}\NormalTok{CARD.INDEX, }\AttributeTok{sep =} \StringTok{"$"}\NormalTok{),}
                        \AttributeTok{CardBrand =}\NormalTok{ cards\_csv}\SpecialCharTok{$}\NormalTok{Card.Brand,}
                        \AttributeTok{CardType =}\NormalTok{ cards\_csv}\SpecialCharTok{$}\NormalTok{Card.Type,}
                        \AttributeTok{CardNumber =}\NormalTok{ cards\_csv}\SpecialCharTok{$}\NormalTok{Card.Number,}
                        \AttributeTok{Expires =}\NormalTok{ cards\_csv}\SpecialCharTok{$}\NormalTok{Expires,}
                        \AttributeTok{CVV =}\NormalTok{ cards\_csv}\SpecialCharTok{$}\NormalTok{CVV, }
                        \AttributeTok{HasChip =}\NormalTok{ cards\_csv}\SpecialCharTok{$}\NormalTok{Has.Chip, }
                        \AttributeTok{CardsIssued =}\NormalTok{ cards\_csv}\SpecialCharTok{$}\NormalTok{Cards.Issued, }
                        \AttributeTok{CreditLimit =}\NormalTok{ cards\_csv}\SpecialCharTok{$}\NormalTok{Credit.Limit,}
                        \AttributeTok{AcctOpenDate =}\NormalTok{ cards\_csv}\SpecialCharTok{$}\NormalTok{Acct.Open.Date,}
                        \AttributeTok{YearPINLastChanged =}\NormalTok{ cards\_csv}\SpecialCharTok{$}\NormalTok{Year.PIN.last.Changed}
\NormalTok{)}
\FunctionTok{head}\NormalTok{(cards, }\DecValTok{3}\NormalTok{)}
\end{Highlighting}
\end{Shaded}

\begin{verbatim}
##   CardId CardBrand CardType   CardNumber Expires CVV HasChip CardsIssued
## 1    0$0      Visa    Debit 4.344677e+15 12/2022 623     YES           2
## 2    0$1      Visa    Debit 4.956966e+15 12/2020 393     YES           2
## 3    0$2      Visa    Debit 4.582313e+15 02/2024 719     YES           2
##   CreditLimit AcctOpenDate YearPINLastChanged
## 1      $24295      09/2002               2008
## 2      $21968      04/2014               2014
## 3      $46414      07/2003               2004
\end{verbatim}

\subsubsection{2.2.d Convert columns - CardBrand, CardType, HasChip to
categorical data
type}\label{d-convert-columns---cardbrand-cardtype-haschip-to-categorical-data-type}

\begin{Shaded}
\begin{Highlighting}[]
\NormalTok{cards}\SpecialCharTok{$}\NormalTok{CardBrand }\OtherTok{\textless{}{-}} \FunctionTok{factor}\NormalTok{(cards}\SpecialCharTok{$}\NormalTok{CardBrand)}
\NormalTok{cards}\SpecialCharTok{$}\NormalTok{CardType }\OtherTok{\textless{}{-}} \FunctionTok{factor}\NormalTok{(cards}\SpecialCharTok{$}\NormalTok{CardType)}
\NormalTok{cards}\SpecialCharTok{$}\NormalTok{HasChip }\OtherTok{\textless{}{-}} \FunctionTok{factor}\NormalTok{(cards}\SpecialCharTok{$}\NormalTok{HasChip)}
\end{Highlighting}
\end{Shaded}

\subsubsection{2.2.e Convert columns - Expires, AcctOpenDate to Date
data
type}\label{e-convert-columns---expires-acctopendate-to-date-data-type}

The Date Format for Expires and AcctOpenDate contains only Month and
Year, but R needs a value for Day as well. The Expires date will be set
to the last day of the month and AcctOpenDate to the first day of the
month for this project.

\begin{Shaded}
\begin{Highlighting}[]
\NormalTok{cards}\SpecialCharTok{$}\NormalTok{Expires }\OtherTok{\textless{}{-}} \FunctionTok{my}\NormalTok{(cards}\SpecialCharTok{$}\NormalTok{Expires)}
\NormalTok{cards}\SpecialCharTok{$}\NormalTok{Expires }\OtherTok{\textless{}{-}} \FunctionTok{ceiling\_date}\NormalTok{(cards}\SpecialCharTok{$}\NormalTok{Expires, }\StringTok{"month"}\NormalTok{) }\SpecialCharTok{{-}} \FunctionTok{days}\NormalTok{(}\DecValTok{1}\NormalTok{)}

\NormalTok{cards}\SpecialCharTok{$}\NormalTok{AcctOpenDate }\OtherTok{\textless{}{-}} \FunctionTok{my}\NormalTok{(cards}\SpecialCharTok{$}\NormalTok{AcctOpenDate)}
\end{Highlighting}
\end{Shaded}

\subsubsection{2.2.f Convert columns - CardNumber, CVV to Text data
type}\label{f-convert-columns---cardnumber-cvv-to-text-data-type}

CardNumber and CVV not used for quantitative statistics, therefore they
are being converted to Text

\begin{Shaded}
\begin{Highlighting}[]
\NormalTok{cards}\SpecialCharTok{$}\NormalTok{CardNumber }\OtherTok{\textless{}{-}} \FunctionTok{as.character}\NormalTok{(cards}\SpecialCharTok{$}\NormalTok{CardNumber)}
\NormalTok{cards}\SpecialCharTok{$}\NormalTok{CVV }\OtherTok{\textless{}{-}} \FunctionTok{as.character}\NormalTok{(cards}\SpecialCharTok{$}\NormalTok{CVV)}
\end{Highlighting}
\end{Shaded}

\subsubsection{2.2.g Convert column CreditLimit to
numeric}\label{g-convert-column-creditlimit-to-numeric}

\begin{Shaded}
\begin{Highlighting}[]
\NormalTok{cards}\SpecialCharTok{$}\NormalTok{CreditLimit }\OtherTok{\textless{}{-}} \FunctionTok{as.numeric}\NormalTok{(}\FunctionTok{sub}\NormalTok{(}\StringTok{\textquotesingle{}.\textquotesingle{}}\NormalTok{, }\StringTok{\textquotesingle{}\textquotesingle{}}\NormalTok{, cards}\SpecialCharTok{$}\NormalTok{CreditLimit))}
\end{Highlighting}
\end{Shaded}

\subsection{2.3 Create table for
merchants}\label{create-table-for-merchants}

\subsubsection{2.3.a Extract Merchant related column from the
Transactions
table}\label{a-extract-merchant-related-column-from-the-transactions-table}

\begin{Shaded}
\begin{Highlighting}[]
\NormalTok{merchants\_temp }\OtherTok{\textless{}{-}} \FunctionTok{data.frame}\NormalTok{(}\AttributeTok{MerchantName =}\NormalTok{ transactions\_csv}\SpecialCharTok{$}\NormalTok{Merchant.Name,}
                            \AttributeTok{MerchantCity =}\NormalTok{ transactions\_csv}\SpecialCharTok{$}\NormalTok{Merchant.City,}
                            \AttributeTok{MerchantState =}\NormalTok{ transactions\_csv}\SpecialCharTok{$}\NormalTok{Merchant.State,}
                            \AttributeTok{MerchantZIP =}\NormalTok{ transactions\_csv}\SpecialCharTok{$}\NormalTok{Zip,}
                            \AttributeTok{MerchantCategoryCode =}\NormalTok{ transactions\_csv}\SpecialCharTok{$}\NormalTok{MCC}
\NormalTok{)}
\end{Highlighting}
\end{Shaded}

\subsubsection{2.3.b Remove redundant data from the Merchants
table}\label{b-remove-redundant-data-from-the-merchants-table}

\begin{Shaded}
\begin{Highlighting}[]
\FunctionTok{nrow}\NormalTok{(merchants\_temp)}
\end{Highlighting}
\end{Shaded}

\begin{verbatim}
## [1] 24386900
\end{verbatim}

\begin{Shaded}
\begin{Highlighting}[]
\FunctionTok{length}\NormalTok{(}\FunctionTok{unique}\NormalTok{(merchants\_temp}\SpecialCharTok{$}\NormalTok{MerchantName))}
\end{Highlighting}
\end{Shaded}

\begin{verbatim}
## [1] 100343
\end{verbatim}

\begin{Shaded}
\begin{Highlighting}[]
\NormalTok{merchants\_temp }\OtherTok{\textless{}{-}}\NormalTok{ merchants\_temp }\SpecialCharTok{\%\textgreater{}\%} \FunctionTok{distinct}\NormalTok{(MerchantName, }\AttributeTok{.keep\_all =} \ConstantTok{TRUE}\NormalTok{)}
\FunctionTok{nrow}\NormalTok{(merchants\_temp)}
\end{Highlighting}
\end{Shaded}

\begin{verbatim}
## [1] 100343
\end{verbatim}

\subsubsection{2.3.c Generate and add MerchantId
column}\label{c-generate-and-add-merchantid-column}

\begin{Shaded}
\begin{Highlighting}[]
\NormalTok{merchants }\OtherTok{\textless{}{-}} \FunctionTok{data.frame}\NormalTok{(}\AttributeTok{MerchantId =} \DecValTok{1}\SpecialCharTok{:}\FunctionTok{nrow}\NormalTok{(merchants\_temp),}
                            \AttributeTok{MerchantName =}\NormalTok{ merchants\_temp}\SpecialCharTok{$}\NormalTok{MerchantName,}
                            \AttributeTok{MerchantCity =}\NormalTok{ merchants\_temp}\SpecialCharTok{$}\NormalTok{MerchantCity,}
                            \AttributeTok{MerchantState =}\NormalTok{ merchants\_temp}\SpecialCharTok{$}\NormalTok{MerchantState,}
                            \AttributeTok{MerchantZIP =}\NormalTok{ merchants\_temp}\SpecialCharTok{$}\NormalTok{MerchantZIP,}
                            \AttributeTok{MerchantCategoryCode =}\NormalTok{ merchants\_temp}\SpecialCharTok{$}\NormalTok{MerchantCategoryCode}
\NormalTok{)}

\FunctionTok{head}\NormalTok{(merchants, }\DecValTok{3}\NormalTok{)}
\end{Highlighting}
\end{Shaded}

\begin{verbatim}
##   MerchantId  MerchantName  MerchantCity MerchantState MerchantZIP
## 1          1  3.527213e+18      La Verne            CA       91750
## 2          2 -7.276121e+17 Monterey Park            CA       91754
## 3          3  3.414527e+18 Monterey Park            CA       91754
##   MerchantCategoryCode
## 1                 5300
## 2                 5411
## 3                 5651
\end{verbatim}

\subsubsection{2.3.d Convert columns - MerchantName, MerchantZIP to Text
data
type}\label{d-convert-columns---merchantname-merchantzip-to-text-data-type}

\begin{Shaded}
\begin{Highlighting}[]
\NormalTok{merchants}\SpecialCharTok{$}\NormalTok{MerchantName }\OtherTok{\textless{}{-}} \FunctionTok{as.character}\NormalTok{(merchants}\SpecialCharTok{$}\NormalTok{MerchantName)}
\NormalTok{merchants}\SpecialCharTok{$}\NormalTok{MerchantZIP }\OtherTok{\textless{}{-}} \FunctionTok{as.character}\NormalTok{(merchants}\SpecialCharTok{$}\NormalTok{MerchantZIP)}
\end{Highlighting}
\end{Shaded}

\subsubsection{2.3.e Convert column MerchantCategoryCode to Categorical
data
type}\label{e-convert-column-merchantcategorycode-to-categorical-data-type}

\begin{Shaded}
\begin{Highlighting}[]
\NormalTok{merchants}\SpecialCharTok{$}\NormalTok{MerchantCategoryCode }\OtherTok{=} \FunctionTok{factor}\NormalTok{(merchants}\SpecialCharTok{$}\NormalTok{MerchantCategoryCode)}
\end{Highlighting}
\end{Shaded}

\subsection{2.4 Create table for date}\label{create-table-for-date}

\subsubsection{2.4.a Transform Time column to categorical (morning,
afternoon, evening,
night)}\label{a-transform-time-column-to-categorical-morning-afternoon-evening-night}

\begin{Shaded}
\begin{Highlighting}[]
\CommentTok{\# Convert column to Time data type}
\NormalTok{transactions\_csv}\SpecialCharTok{$}\NormalTok{Time }\OtherTok{\textless{}{-}} \FunctionTok{as.POSIXct}\NormalTok{(}\FunctionTok{strptime}\NormalTok{(transactions\_csv}\SpecialCharTok{$}\NormalTok{Time,}\StringTok{"\%H:\%M"}\NormalTok{),}\StringTok{"UTC"}\NormalTok{)}

\CommentTok{\# Define time range for each category}
\NormalTok{x}\OtherTok{=}\FunctionTok{as.POSIXct}\NormalTok{(}\FunctionTok{strptime}\NormalTok{(}\FunctionTok{c}\NormalTok{(}\StringTok{"050000"}\NormalTok{,}\StringTok{"105959"}\NormalTok{,}\StringTok{"110000"}\NormalTok{,}\StringTok{"155959"}\NormalTok{,}\StringTok{"160000"}\NormalTok{,}
                        \StringTok{"185959"}\NormalTok{),}\StringTok{"\%H\%M\%S"}\NormalTok{),}\StringTok{"UTC"}\NormalTok{)}

\CommentTok{\# Transform Time column }
\NormalTok{transactions\_csv}\SpecialCharTok{$}\NormalTok{Time }\OtherTok{=} \FunctionTok{case\_when}\NormalTok{(}
  \FunctionTok{between}\NormalTok{(transactions\_csv}\SpecialCharTok{$}\NormalTok{Time,x[}\DecValTok{1}\NormalTok{],x[}\DecValTok{2}\NormalTok{]) }\SpecialCharTok{\textasciitilde{}}\StringTok{"morning"}\NormalTok{,}
  \FunctionTok{between}\NormalTok{(transactions\_csv}\SpecialCharTok{$}\NormalTok{Time,x[}\DecValTok{3}\NormalTok{],x[}\DecValTok{4}\NormalTok{]) }\SpecialCharTok{\textasciitilde{}}\StringTok{"afternoon"}\NormalTok{,}
  \FunctionTok{between}\NormalTok{(transactions\_csv}\SpecialCharTok{$}\NormalTok{Time,x[}\DecValTok{5}\NormalTok{],x[}\DecValTok{6}\NormalTok{]) }\SpecialCharTok{\textasciitilde{}}\StringTok{"evening"}\NormalTok{,}
  \ConstantTok{TRUE} \SpecialCharTok{\textasciitilde{}}\StringTok{"night"}\NormalTok{)}

\FunctionTok{head}\NormalTok{(transactions\_csv}\SpecialCharTok{$}\NormalTok{Time)}
\end{Highlighting}
\end{Shaded}

\begin{verbatim}
## [1] "morning"   "morning"   "morning"   "evening"   "morning"   "afternoon"
\end{verbatim}

\subsubsection{2.4.b Extract Date related column from the Transactions
table}\label{b-extract-date-related-column-from-the-transactions-table}

\begin{Shaded}
\begin{Highlighting}[]
\NormalTok{date\_temp }\OtherTok{\textless{}{-}} \FunctionTok{data.frame}\NormalTok{(}\AttributeTok{DateYear =}\NormalTok{ transactions\_csv}\SpecialCharTok{$}\NormalTok{Year,}
                        \AttributeTok{DateMonth =}\NormalTok{ transactions\_csv}\SpecialCharTok{$}\NormalTok{Month,}
                        \AttributeTok{DateDay =}\NormalTok{ transactions\_csv}\SpecialCharTok{$}\NormalTok{Day,}
                        \AttributeTok{Time =}\NormalTok{ transactions\_csv}\SpecialCharTok{$}\NormalTok{Time}
\NormalTok{)}
\end{Highlighting}
\end{Shaded}

\subsubsection{2.4.c Remove redundant data from the Date
table}\label{c-remove-redundant-data-from-the-date-table}

\begin{Shaded}
\begin{Highlighting}[]
\NormalTok{date\_temp }\OtherTok{\textless{}{-}}\NormalTok{ date\_temp }\SpecialCharTok{\%\textgreater{}\%} \FunctionTok{distinct}\NormalTok{() }
\end{Highlighting}
\end{Shaded}

\subsubsection{2.4.d Generate and add DateId
column}\label{d-generate-and-add-dateid-column}

\begin{Shaded}
\begin{Highlighting}[]
\NormalTok{date }\OtherTok{\textless{}{-}} \FunctionTok{data.frame}\NormalTok{(}\AttributeTok{DateId =} \DecValTok{1}\SpecialCharTok{:}\FunctionTok{nrow}\NormalTok{(date\_temp),}
                        \AttributeTok{DateYear =}\NormalTok{ date\_temp}\SpecialCharTok{$}\NormalTok{DateYear,}
                        \AttributeTok{DateMonth =}\NormalTok{ date\_temp}\SpecialCharTok{$}\NormalTok{DateMonth,}
                        \AttributeTok{DateDay =}\NormalTok{ date\_temp}\SpecialCharTok{$}\NormalTok{DateDay,}
                        \AttributeTok{Time =}\NormalTok{ date\_temp}\SpecialCharTok{$}\NormalTok{Time}
\NormalTok{)}
\end{Highlighting}
\end{Shaded}

\subsubsection{2.4.e Convert column Time to Categorical data
type}\label{e-convert-column-time-to-categorical-data-type}

\begin{Shaded}
\begin{Highlighting}[]
\NormalTok{date}\SpecialCharTok{$}\NormalTok{Time }\OtherTok{=} \FunctionTok{factor}\NormalTok{(date}\SpecialCharTok{$}\NormalTok{Time)}
\end{Highlighting}
\end{Shaded}

\subsection{2.5 Fact table -
transactions}\label{fact-table---transactions}

\subsubsection{2.4.a Replace existing merchant columns with single
merchant
id}\label{a-replace-existing-merchant-columns-with-single-merchant-id}

\begin{Shaded}
\begin{Highlighting}[]
\CommentTok{\# Uses Merchant Name (unique column in merchants table) }
\CommentTok{\# to assign Merchant Id values in Transactions table}
\NormalTok{transactions\_csv}\SpecialCharTok{$}\NormalTok{MerchantId }\OtherTok{\textless{}{-}}\NormalTok{ transactions\_csv}\SpecialCharTok{$}\NormalTok{Merchant.Name}
\NormalTok{conditions }\OtherTok{\textless{}{-}}\NormalTok{ merchants}\SpecialCharTok{$}\NormalTok{MerchantName}
\NormalTok{replacement\_values }\OtherTok{\textless{}{-}}\NormalTok{ merchants}\SpecialCharTok{$}\NormalTok{MerchantId}
\NormalTok{transactions\_csv}\SpecialCharTok{$}\NormalTok{MerchantId }\OtherTok{\textless{}{-}} \FunctionTok{replace}\NormalTok{(transactions\_csv}\SpecialCharTok{$}\NormalTok{MerchantId, }
\NormalTok{                                   transactions\_csv}\SpecialCharTok{$}\NormalTok{MerchantId }\SpecialCharTok{\%in\%}\NormalTok{ conditions, }
\NormalTok{                                   replacement\_values)}
\end{Highlighting}
\end{Shaded}

\subsubsection{2.4.b Replace existing date columns with single date
id}\label{b-replace-existing-date-columns-with-single-date-id}

\begin{Shaded}
\begin{Highlighting}[]
\CommentTok{\# Combine all columns in Dates table}
\NormalTok{temp\_dates }\OtherTok{=} \FunctionTok{data.frame}\NormalTok{(}\AttributeTok{DateId =}\NormalTok{ date}\SpecialCharTok{$}\NormalTok{DateId,}
                        \AttributeTok{Condition =} \FunctionTok{paste}\NormalTok{(date}\SpecialCharTok{$}\NormalTok{DateYear,}
\NormalTok{                                          date}\SpecialCharTok{$}\NormalTok{DateMonth,}
\NormalTok{                                          date}\SpecialCharTok{$}\NormalTok{DateDay,}
\NormalTok{                                          date}\SpecialCharTok{$}\NormalTok{Time)}
\NormalTok{)}

\CommentTok{\# Combine all columns that were added to Date table from Transaction table }
\NormalTok{transactions\_csv}\SpecialCharTok{$}\NormalTok{Condition }\OtherTok{\textless{}{-}} \FunctionTok{paste}\NormalTok{(transactions\_csv}\SpecialCharTok{$}\NormalTok{Year,}
\NormalTok{                                transactions\_csv}\SpecialCharTok{$}\NormalTok{Month,}
\NormalTok{                                transactions\_csv}\SpecialCharTok{$}\NormalTok{Day,}
\NormalTok{                                transactions\_csv}\SpecialCharTok{$}\NormalTok{Time)}

\CommentTok{\# Add DateId column with temporary values}
\NormalTok{transactions\_csv}\SpecialCharTok{$}\NormalTok{DateId }\OtherTok{\textless{}{-}} \DecValTok{1}\SpecialCharTok{:}\FunctionTok{nrow}\NormalTok{(transactions\_csv)}

\CommentTok{\# Replace DateId based on the combined Date columns in both tables}
\NormalTok{transactions\_csv}\SpecialCharTok{$}\NormalTok{DateId }\OtherTok{\textless{}{-}} \FunctionTok{replace}\NormalTok{(transactions\_csv}\SpecialCharTok{$}\NormalTok{DateID, }
\NormalTok{                             transactions\_csv}\SpecialCharTok{$}\NormalTok{Condition }\SpecialCharTok{\%in\%}\NormalTok{ temp\_dates}\SpecialCharTok{$}\NormalTok{Condition, }
\NormalTok{                             temp\_dates}\SpecialCharTok{$}\NormalTok{DateId)}
\end{Highlighting}
\end{Shaded}

\subsubsection{2.4.c Remove duplicates}\label{c-remove-duplicates}

\begin{Shaded}
\begin{Highlighting}[]
\NormalTok{transactions\_csv }\OtherTok{\textless{}{-}}\NormalTok{ transactions\_csv }\SpecialCharTok{\%\textgreater{}\%} \FunctionTok{distinct}\NormalTok{() }
\end{Highlighting}
\end{Shaded}

\subsubsection{2.4.d Combine user id and card index to generate CardId
column}\label{d-combine-user-id-and-card-index-to-generate-cardid-column}

\subsubsection{2.4.e Rename column names}\label{e-rename-column-names}

\begin{Shaded}
\begin{Highlighting}[]
\NormalTok{transactions }\OtherTok{\textless{}{-}} \FunctionTok{data.frame}\NormalTok{(}\AttributeTok{TransactionId =} \DecValTok{1}\SpecialCharTok{:}\FunctionTok{nrow}\NormalTok{(transactions\_csv),}
                              \AttributeTok{CardId =} \FunctionTok{paste}\NormalTok{(transactions\_csv}\SpecialCharTok{$}\NormalTok{User, }
\NormalTok{                                             transactions\_csv}\SpecialCharTok{$}\NormalTok{Card, }\AttributeTok{sep =} \StringTok{"$"}\NormalTok{),}
                              \AttributeTok{UserId =}\NormalTok{ transactions\_csv}\SpecialCharTok{$}\NormalTok{User,}
                              \AttributeTok{DateId =}\NormalTok{ transactions\_csv}\SpecialCharTok{$}\NormalTok{DateId,}
                              \AttributeTok{MerchantId =}\NormalTok{ transactions\_csv}\SpecialCharTok{$}\NormalTok{MerchantId,}
                              \AttributeTok{Amount =}\NormalTok{ transactions\_csv}\SpecialCharTok{$}\NormalTok{Amount, }
                              \AttributeTok{UseChip =}\NormalTok{ transactions\_csv}\SpecialCharTok{$}\NormalTok{Use.Chip, }
                              \AttributeTok{Errors =}\NormalTok{ transactions\_csv}\SpecialCharTok{$}\NormalTok{Errors., }
                              \AttributeTok{IsFraud =}\NormalTok{ transactions\_csv}\SpecialCharTok{$}\NormalTok{Is.Fraud}
\NormalTok{)}

\FunctionTok{head}\NormalTok{(transactions, }\DecValTok{3}\NormalTok{)}
\end{Highlighting}
\end{Shaded}

\begin{verbatim}
##   TransactionId CardId UserId DateId MerchantId  Amount           UseChip
## 1             1    0$0      0      1          1 $134.09 Swipe Transaction
## 2             2    0$0      0      2          2  $38.48 Swipe Transaction
## 3             3    0$0      0      3          3 $120.34 Swipe Transaction
##   Errors IsFraud
## 1             No
## 2             No
## 3             No
\end{verbatim}

\subsubsection{2.4.f Covert columns - UseChip, IsFraud to categorical
data
type}\label{f-covert-columns---usechip-isfraud-to-categorical-data-type}

\begin{Shaded}
\begin{Highlighting}[]
\NormalTok{transactions}\SpecialCharTok{$}\NormalTok{UseChip }\OtherTok{\textless{}{-}} \FunctionTok{factor}\NormalTok{(transactions}\SpecialCharTok{$}\NormalTok{UseChip)}
\NormalTok{transactions}\SpecialCharTok{$}\NormalTok{IsFraud }\OtherTok{\textless{}{-}} \FunctionTok{factor}\NormalTok{(transactions}\SpecialCharTok{$}\NormalTok{IsFraud)}
\end{Highlighting}
\end{Shaded}

\subsubsection{2.2.g Convert column Amount to
numeric}\label{g-convert-column-amount-to-numeric}

\begin{Shaded}
\begin{Highlighting}[]
\NormalTok{transactions}\SpecialCharTok{$}\NormalTok{Amount }\OtherTok{\textless{}{-}} \FunctionTok{as.numeric}\NormalTok{(}\FunctionTok{sub}\NormalTok{(}\StringTok{\textquotesingle{}.\textquotesingle{}}\NormalTok{, }\StringTok{\textquotesingle{}\textquotesingle{}}\NormalTok{, transactions}\SpecialCharTok{$}\NormalTok{Amount))}
\end{Highlighting}
\end{Shaded}

\section{3. LOAD}\label{load}

\subsection{Dimension table -
users\_dim}\label{dimension-table---users_dim}

\begin{Shaded}
\begin{Highlighting}[]
\NormalTok{users\_dim }\OtherTok{\textless{}{-}}\NormalTok{ users }\SpecialCharTok{\%\textgreater{}\%} \FunctionTok{select}\NormalTok{(UserId, FullName, CurrentAge, RetirementAge, }
\NormalTok{                              BirthYear, BirthMonth, Gender, Address1, Address2, }
\NormalTok{                              City, State, Zipcode, Latitude, Longitude, }
\NormalTok{                              PerCapitaIncomeZipcode, YearlyIncome, TotalDebt, }
\NormalTok{                              FICOScore, NumCreditCards) }\SpecialCharTok{\%\textgreater{}\%} \FunctionTok{distinct}\NormalTok{()}
\FunctionTok{summary}\NormalTok{(users\_dim)}
\end{Highlighting}
\end{Shaded}

\begin{verbatim}
##      UserId         FullName           CurrentAge     RetirementAge  
##  Min.   :   1.0   Length:2000        Min.   : 18.00   Min.   :50.00  
##  1st Qu.: 500.8   Class :character   1st Qu.: 30.00   1st Qu.:65.00  
##  Median :1000.5   Mode  :character   Median : 44.00   Median :66.00  
##  Mean   :1000.5                      Mean   : 45.39   Mean   :66.24  
##  3rd Qu.:1500.2                      3rd Qu.: 58.00   3rd Qu.:68.00  
##  Max.   :2000.0                      Max.   :101.00   Max.   :79.00  
##    BirthYear      BirthMonth        Gender       Address1        
##  Min.   :1918   Min.   : 1.000   Female:1016   Length:2000       
##  1st Qu.:1961   1st Qu.: 3.000   Male  : 984   Class :character  
##  Median :1975   Median : 7.000                 Mode  :character  
##  Mean   :1974   Mean   : 6.439                                   
##  3rd Qu.:1989   3rd Qu.:10.000                                   
##  Max.   :2002   Max.   :12.000                                   
##    Address2             City              State             Zipcode         
##  Length:2000        Length:2000        Length:2000        Length:2000       
##  Class :character   Class :character   Class :character   Class :character  
##  Mode  :character   Mode  :character   Mode  :character   Mode  :character  
##                                                                             
##                                                                             
##                                                                             
##     Latitude       Longitude       PerCapitaIncomeZipcode  YearlyIncome   
##  Min.   :20.88   Min.   :-159.41   Min.   :     0         Min.   :     1  
##  1st Qu.:33.84   1st Qu.: -97.39   1st Qu.: 16824         1st Qu.: 32818  
##  Median :38.25   Median : -86.44   Median : 20581         Median : 40744  
##  Mean   :37.39   Mean   : -91.55   Mean   : 23142         Mean   : 45716  
##  3rd Qu.:41.20   3rd Qu.: -80.13   3rd Qu.: 26286         3rd Qu.: 52698  
##  Max.   :61.20   Max.   : -68.67   Max.   :163145         Max.   :307018  
##    TotalDebt        FICOScore     NumCreditCards 
##  Min.   :     0   Min.   :480.0   Min.   :1.000  
##  1st Qu.: 23987   1st Qu.:681.0   1st Qu.:2.000  
##  Median : 58251   Median :711.5   Median :3.000  
##  Mean   : 63710   Mean   :709.7   Mean   :3.073  
##  3rd Qu.: 89070   3rd Qu.:753.0   3rd Qu.:4.000  
##  Max.   :516263   Max.   :850.0   Max.   :9.000
\end{verbatim}

\subsection{Dimension table -
cards\_dim}\label{dimension-table---cards_dim}

\begin{Shaded}
\begin{Highlighting}[]
\NormalTok{cards\_dim }\OtherTok{\textless{}{-}}\NormalTok{ cards }\SpecialCharTok{\%\textgreater{}\%} \FunctionTok{select}\NormalTok{(CardId, CardBrand, CardType, CardNumber, Expires, }
\NormalTok{                              CVV, HasChip, CardsIssued, CreditLimit, }
\NormalTok{                              AcctOpenDate, YearPINLastChanged) }\SpecialCharTok{\%\textgreater{}\%} \FunctionTok{distinct}\NormalTok{()}
\FunctionTok{summary}\NormalTok{(cards)}
\end{Highlighting}
\end{Shaded}

\begin{verbatim}
##     CardId               CardBrand               CardType     CardNumber       
##  Length:6146        Amex      : 402   Credit         :2057   Length:6146       
##  Class :character   Discover  : 209   Debit          :3511   Class :character  
##  Mode  :character   Mastercard:3209   Debit (Prepaid): 578   Mode  :character  
##                     Visa      :2326                                            
##                                                                                
##                                                                                
##     Expires               CVV            HasChip     CardsIssued   
##  Min.   :1997-07-31   Length:6146        NO : 646   Min.   :1.000  
##  1st Qu.:2020-02-29   Class :character   YES:5500   1st Qu.:1.000  
##  Median :2021-09-30   Mode  :character              Median :1.000  
##  Mean   :2020-11-06                                 Mean   :1.503  
##  3rd Qu.:2023-05-31                                 3rd Qu.:2.000  
##  Max.   :2024-12-31                                 Max.   :3.000  
##   CreditLimit      AcctOpenDate        YearPINLastChanged
##  Min.   :     0   Min.   :1991-01-01   Min.   :2002      
##  1st Qu.:  7043   1st Qu.:2006-10-01   1st Qu.:2010      
##  Median : 12592   Median :2010-02-15   Median :2013      
##  Mean   : 14347   Mean   :2011-01-15   Mean   :2013      
##  3rd Qu.: 19156   3rd Qu.:2016-05-01   3rd Qu.:2017      
##  Max.   :151223   Max.   :2020-02-01   Max.   :2020
\end{verbatim}

\subsection{Dimension table -
cards\_dim}\label{dimension-table---cards_dim-1}

\begin{Shaded}
\begin{Highlighting}[]
\NormalTok{merchants\_dim }\OtherTok{\textless{}{-}}\NormalTok{ merchants }\SpecialCharTok{\%\textgreater{}\%} \FunctionTok{select}\NormalTok{(MerchantId, MerchantName, MerchantCity, }
\NormalTok{                                      MerchantState, MerchantZIP, }
\NormalTok{                                      MerchantCategoryCode) }\SpecialCharTok{\%\textgreater{}\%} \FunctionTok{distinct}\NormalTok{()}
\FunctionTok{summary}\NormalTok{(merchants\_dim)}
\end{Highlighting}
\end{Shaded}

\begin{verbatim}
##    MerchantId     MerchantName       MerchantCity       MerchantState     
##  Min.   :     1   Length:100343      Length:100343      Length:100343     
##  1st Qu.: 25086   Class :character   Class :character   Class :character  
##  Median : 50172   Mode  :character   Mode  :character   Mode  :character  
##  Mean   : 50172                                                           
##  3rd Qu.: 75258                                                           
##  Max.   :100343                                                           
##                                                                           
##  MerchantZIP        MerchantCategoryCode
##  Length:100343      5411   :10051       
##  Class :character   5812   : 7910       
##  Mode  :character   5912   : 6150       
##                     5813   : 5134       
##                     5921   : 4701       
##                     4900   : 4695       
##                     (Other):61702
\end{verbatim}

\subsection{Dimension table -
date\_dim}\label{dimension-table---date_dim}

\begin{Shaded}
\begin{Highlighting}[]
\NormalTok{date\_dim }\OtherTok{\textless{}{-}}\NormalTok{ date }\SpecialCharTok{\%\textgreater{}\%} \FunctionTok{select}\NormalTok{(DateId, DateYear, DateMonth, }
\NormalTok{                            DateDay, Time) }\SpecialCharTok{\%\textgreater{}\%} \FunctionTok{distinct}\NormalTok{()}
\FunctionTok{summary}\NormalTok{(date\_dim)}
\end{Highlighting}
\end{Shaded}

\begin{verbatim}
##      DateId         DateYear      DateMonth         DateDay     
##  Min.   :    1   Min.   :1991   Min.   : 1.000   Min.   : 1.00  
##  1st Qu.:10366   1st Qu.:1998   1st Qu.: 3.000   1st Qu.: 8.00  
##  Median :20731   Median :2005   Median : 7.000   Median :16.00  
##  Mean   :20731   Mean   :2005   Mean   : 6.514   Mean   :15.73  
##  3rd Qu.:31096   3rd Qu.:2013   3rd Qu.:10.000   3rd Qu.:23.00  
##  Max.   :41461   Max.   :2020   Max.   :12.000   Max.   :31.00  
##         Time      
##  afternoon:10554  
##  evening  :10374  
##  morning  :10431  
##  night    :10102  
##                   
## 
\end{verbatim}

\subsection{Fact table -
transactions\_fact}\label{fact-table---transactions_fact}

\begin{Shaded}
\begin{Highlighting}[]
\NormalTok{transactions\_fact }\OtherTok{\textless{}{-}}\NormalTok{ transactions }\SpecialCharTok{\%\textgreater{}\%} \FunctionTok{select}\NormalTok{(TransactionId, UserId, CardId, }
\NormalTok{                                             DateId, Amount, UseChip, Errors, }
\NormalTok{                                             IsFraud) }\SpecialCharTok{\%\textgreater{}\%} \FunctionTok{distinct}\NormalTok{()}
\FunctionTok{summary}\NormalTok{(transactions\_fact)}
\end{Highlighting}
\end{Shaded}

\begin{verbatim}
##  TransactionId          UserId        CardId              DateId     
##  Min.   :       1   Min.   :   0   Length:24386900    Min.   :    1  
##  1st Qu.: 6096726   1st Qu.: 510   Class :character   1st Qu.:10356  
##  Median :12193450   Median :1006   Mode  :character   Median :20724  
##  Mean   :12193450   Mean   :1001                      Mean   :20726  
##  3rd Qu.:18290175   3rd Qu.:1477                      3rd Qu.:31093  
##  Max.   :24386900   Max.   :1999                      Max.   :41461  
##      Amount                       UseChip            Errors         
##  Min.   : -500.00   Chip Transaction  : 6287598   Length:24386900   
##  1st Qu.:    9.20   Online Transaction: 2713220   Class :character  
##  Median :   30.14   Swipe Transaction :15386082   Mode  :character  
##  Mean   :   43.63                                                   
##  3rd Qu.:   65.06                                                   
##  Max.   :12390.50                                                   
##  IsFraud       
##  No :24357143  
##  Yes:   29757  
##                
##                
##                
## 
\end{verbatim}

\end{document}
